% 原PDF158页附录4.E标题—160页末。
\section*{附录4.E\quad 动态结构方程模型参数的最大似然估计}
\addcontentsline{toc}{section}{附录4.E 动态结构方程模型参数的最大似然估计}
\markright{附录4.E 动态结构方程模型参数的最大似然估计}

令 $(y_i^t)'$、$(y_i^{(t-1)}(x_i)'$ 和 $(e_i^t)'$，分别为矩阵 $Y^t,Y^{t-1},X,E^{(t)}$ 的第 $i$ 行。换言之，定义
\begin{equation}
\begin{aligned}
(y_i^t)'&=\begin{bmatrix}y_{i1}^t&y_{i2}^t&\cdots&y_{iM}^t\end{bmatrix},\\
(y_{li}^{t-1})'&=\begin{bmatrix}y_{li1}^{t-1}&y_{li2}^{t-1}&\cdots&y_{liM}^{t-1}\end{bmatrix},\\
x_i'&=\begin{bmatrix}x_{i1}&x_{i2}&\cdots&x_{iK}\end{bmatrix},\\
(e_i^t)'&=\begin{bmatrix}e_{i1}^t&e_{i2}^t&\cdots&e_{iM}^t\end{bmatrix},
\quad i=1,\ldots,n.
\end{aligned}\tag{4.E.1}
\end{equation}
由式（4.173），
\begin{equation}
(y_i^t)'\Gamma^t+(y_{li}^{t-1})'A^t+x_i'B^t+(e_i^t)'=0.\tag{4.E.2}
\end{equation}
令
\[
e_i^t=\begin{bmatrix}e_{i1}^t\\\vdots\\e_{iM}^t\end{bmatrix},
\quad t=1,\ldots,T,\qquad
e_i=\begin{bmatrix}e_i^1\\\vdots\\e_i^T\end{bmatrix},\quad i=1,\ldots,n.
\]
误差变量之间的协方差为
\begin{equation}
\operatorname{cov}(e_i,e_i)=\Sigma=\begin{bmatrix}
\Sigma_{11}&\Sigma_{12}&\cdots&\Sigma_{1T}\\
\Sigma_{21}&\Sigma_{22}&\cdots&\Sigma_{2T}\\
\vdots&\vdots&\vdots&\vdots\\
\Sigma_{T1}&\Sigma_{T12}&\cdots&\Sigma_{TT}
\end{bmatrix},\quad i=1,\ldots,n,\tag{4.E.3}
\end{equation}
其中
\begin{equation}
\Sigma_{tt}=\operatorname{cov}(e_i^t,e_i^t)=\begin{bmatrix}
\sigma_{11}^{tt}&\cdots&\sigma_{1M}^{tt}\\
\vdots&\vdots&\vdots\\\sigma_{M1}^{tt}&\cdots&\sigma_{MM}^{tt}
\end{bmatrix}.\tag{4.E.4}
\end{equation}
假设 $\Gamma^t$ 非奇异。由式（4.E.2），向量 $(y_i^t)'$ 可以写成误差 $(e_i^t)'$ 的函数：
\begin{equation}
\begin{aligned}
(y_i^t)'={}&-(e_i^t)'(\Gamma^t)^{-1}
-(y_{li}^{t-1})'A^t(\Gamma^t)^{-1}
-x_i'B^t(\Gamma^t)^{-1},\\
&i=1,\ldots,n,\quad t=1,\ldots,T
\end{aligned}\tag{4.E.5}
\end{equation}
或
\begin{equation}
\begin{aligned}
y_i^t={}&-((\Gamma^t)')^{-1}e_i^t
-((\Gamma^t)')^{-1}(A^t)'y_{li}^{t-1}
-((\Gamma^t)')^{-1}(B^t)'x_i,\\
&i=1,\ldots,n,\quad t=1,\ldots,T.
\end{aligned}\tag{4.E.6}
\end{equation}
令
\begin{equation}
y_i=\begin{bmatrix}y_i^1\\y_i^2\\\vdots\\y_i^T\end{bmatrix},\quad
y_{li}=\begin{bmatrix}0\\y_{li}^1\\\vdots\\y_{li}^{T-1}\end{bmatrix},\quad
e_i=\begin{bmatrix}e_i^1\\e_i^2\\\vdots\\e_i^T\end{bmatrix},\quad i=1,\ldots,n.\tag{4.E.7}
\end{equation}
假设误差向量服从均值为零、协方差矩阵为 $\Sigma$ 的多元正态分布，则
\begin{equation}
g(e_i)=(2\pi)^{-TM/2}|\Sigma|^{-1/2}
\exp\left\{-\frac12e_i'\Sigma^{-1}e_i\right\}.\tag{4.E.8}
\end{equation}
基因表达向量 $y_i^t$ 是误差的函数（式（4.E.6）），其Jacobian矩阵为
\begin{equation}
\begin{aligned}
J=\frac{\partial y_i}{\partial e_i'}
&=\begin{bmatrix}
\dfrac{\partial y_i^1}{\partial e_i'}\\
\dfrac{\partial y_i^2}{\partial e_i'}\\\vdots\\
\dfrac{\partial y_i^T}{\partial e_i'}
\end{bmatrix}\\
&=\begin{bmatrix}
-((\Gamma^1)')^{-1}&0&\cdots&0\\
0&-((\Gamma^2)')^{-1}&\cdots&0\\
\vdots&\vdots&\vdots&\vdots\\
0&0&\cdots&-((\Gamma^T)')^{-1}
\end{bmatrix},\quad i=1,\ldots,n.
\end{aligned}\tag{4.E.9}
\end{equation}
于是，$y_i$ 的密度函数为
\begin{equation}
\begin{aligned}
f(y_i)&=|J|^{-1}g(e_i)\\
&=\prod_{t=1}^{T}|\Gamma^t|g(e_i)\\
&=(2\pi)^{-TM/2}|\Sigma|^{-1/2}\prod_{t=1}^{T}|\Gamma^t|
\exp\left\{-\frac12e_i'\Sigma^{-1}e_i\right\},
\end{aligned}\tag{4.E.10}
\end{equation}
其中
\[
e_i^t=-(\Gamma^t)'y_i^t-(A^t)'y_{li}^{t-1}-(B^t)'X_i,
\quad t=1,\ldots,T,\qquad
e_i=\begin{bmatrix}e_i^1\\e_i^2\\\vdots\\e_i^T\end{bmatrix}.
\]
因此，基因表达的联合概率为
\begin{equation}
f(y_1,\ldots,y_n)=(2\pi)^{-nTM/2}|\Sigma|^{-n/2}
\prod_{t=1}^{T}|\Gamma^t|^n
\exp\left\{-\frac12\sum_{i=1}^{n}e_i'\Sigma^{-1}e_i\right\}.\tag{4.E.11}
\end{equation}
对数似然函数定义为
\begin{equation}
\begin{aligned}
l(\Gamma,A,B,\Sigma\mid Y)
={}&-\frac{nTM}{2}\log(2\pi)-\frac n2\log|\Sigma|
+n\sum_{t=1}^{T}\log|\Gamma^t|\\
&{}-\frac12\sum_{i=1}^{n}e_i'\Sigma^{-1}e_i.
\end{aligned}\tag{4.E.12}
\end{equation}
若假设误差 $e_i^t$ 与 $e_i^t$ 不相关，则 $\Sigma$ 变成分块对角矩阵：
\[
\Sigma=\begin{bmatrix}
\Sigma_{11}&0&\cdots&0\\0&\Sigma_{22}&\cdots&0\\
\vdots&\vdots&\vdots&\vdots\\0&0&\cdots&\Sigma_{TT}
\end{bmatrix}.
\]
于是，似然函数化为
\begin{equation}
\begin{aligned}
l(\Gamma,A,B,\Sigma\mid Y)
={}&-\frac{nTM}{2}\log(2\pi)-\frac n2\sum_{t=1}^{T}\log|\Sigma_{tt}|
+n\sum_{t=1}^{T}\log|\Gamma^t|\\
&{}-\frac12\sum_{i=1}^{n}\sum_{t=1}^{T}(e_i^t)'\Sigma_{tt}^{-1}e_i^t.
\end{aligned}\tag{4.E.13}
\end{equation}
不同时间点误差相关时，极大化式（4.E.12）；不相关时，极大化式（4.E.13），即可得到结构方程参数的估计量 $\widehat\Gamma,\widehat A,\widehat B,\widehat\Sigma$。
